\plannedchapter{uops、执行端口和 llvm-mca}
{本章讲 x86 指令如何被解码为 micro-ops，以及如何用 \code{llvm-mca} 建立静态后端模型。}
{\item 一条汇编指令为什么可能是多个 uops？
\item port pressure 如何限制吞吐？
\item \code{llvm-mca} 能解释什么，不能解释什么？}
{\item uop decomposition、execution resources、load/store ports。
\item block throughput、resource pressure、timeline。
\item 静态模型与实测差异。}
{抽取 SAXPY/FMA loop，用 \code{llvm-mca} 分析 uops、resource pressure 和 block throughput。}
